\documentclass[11pt]{article}
\usepackage[T1]{fontenc}
\usepackage[utf8]{inputenc}
\usepackage[english]{babel}
\usepackage[margin=1in]{geometry}
\usepackage{amsmath,amssymb}
\usepackage{enumitem}
\usepackage{microtype}

\ifdefined\pdfcatalog
  \pdfcatalog{/Lang (en-US)}
\fi

\setlength{\parindent}{0pt}
\setlength{\parskip}{0.55em}
\setlength{\emergencystretch}{2em}
\setlist{nosep,leftmargin=*}

\newcommand{\findinglabel}[1]{\textbf{#1}}
\newcommand{\classification}[1]{\textit{Classification: #1.}}

\title{Technical Review of ``Early-Stopping Activation Steering for Hallucination Mitigation in Vietnamese Clinical Language Models''}
\author{Manuscript review}
\date{August 15, 2026}

\begin{document}
\selectlanguage{english}
\maketitle

\section*{Overall assessment}

The revised manuscript now evaluates a genuine finite-$K$ intervention and reports uncertainty for part of the sampling experiment.  These are useful improvements.  However, the current experiments still do not isolate the effect of early stopping, because the early-stop and full-steering conditions use different initial scales and radically different cumulative intervention magnitudes.  The paper also evaluates lexical/embedding similarity and repetition rather than clinical hallucination correctness, appears to select configurations on the reported test set, and contains ambiguous or internally incorrect statistical and quantitative claims.  The source does not compile as written.  The paper therefore requires major methodological revision before its title-level claim of hallucination mitigation is supported.

\section*{Major technical findings}

\subsection*{1. The early-stopping comparison is confounded by steering strength}

\classification{Definite experimental-design confound}

\findinglabel{Location.} Methodology, Equations~(2)--(3); sampling and greedy result tables; Key Finding~1; Conclusion.

\findinglabel{Problem.} Early stopping changes both the temporal schedule and the steering magnitude.  The sampling comparisons use $(K=16,\alpha_0=15)$ versus full steering with $(K=\infty,\alpha=20)$ at 80 tokens, and $(K=16,\alpha_0=18)$ versus $(K=\infty,\alpha=20)$ at 200 tokens.  For the stated decay,
\[
\sum_{t=1}^{16}\alpha(t)
=\alpha_0\sum_{j=0}^{15}\left(1-\frac{j}{16}\right)
=8.5\alpha_0.
\]
Thus the cumulative early-stop intervention is $127.5$ for $\alpha_0=15$ and $153$ for $\alpha_0=18$, whereas full steering at $\alpha=20$ has cumulative magnitudes $1600$ over 80 tokens and $4000$ over 200 tokens.  Reduced repetition can therefore be caused by the much weaker total perturbation rather than by the stopping/decay schedule itself.

\findinglabel{Why it matters.} This comparison cannot support the causal statement that early stopping suppresses ``representation saturation'' or repetition.  It only shows that one weaker, shorter schedule differs from one stronger, continuous schedule.

\findinglabel{Suggested fix.} Add a factorial ablation that holds $\alpha_0$ fixed while varying the schedule: continuous steering, hard cutoff at $K=16$, and linear decay at $K=16$.  Also include controls matched for cumulative intervention magnitude so that schedule shape and total dose can be separated.  Report results over several $K$ and $\alpha_0$ values selected exclusively on a validation set.  Until then, describe the finding as an association for the tested configurations, not an effect caused specifically by early stopping.

\subsection*{2. The paper still does not measure clinical hallucination mitigation}

\classification{Unsupported central claim}

\findinglabel{Location.} Title; Abstract; Introduction; benchmark description; result tables; Conclusion.

\findinglabel{Problem.} ROUGE-L, BERTScore, and 4-gram repetition are the only reported endpoints.  BERTScore is a semantic-similarity measure, but it does not directly determine whether a dosage, pregnancy contraindication, or drug-interaction mechanism is clinically correct.  Repetition is a fluency artifact, not a hallucination label.  No expert correctness/safety evaluation or category-specific clinical error rate is reported.

\findinglabel{Why it matters.} A response can be lexically and semantically similar to a reference while reversing a dosage or contraindication.  Conversely, a correct paraphrase may receive a lower overlap score.  Reduced repetition relative to full steering does not establish safer medical content.

\findinglabel{Suggested fix.} Make blinded clinical correctness the primary endpoint.  Report expert-adjudicated accuracy and unsafe-error rates separately for special dosage, pregnancy/lactation, and interaction cases, with paired confidence intervals.  Add contradiction/entailment evaluation validated for Vietnamese medical text and retain ROUGE-L/BERTScore as secondary measures.  If this evaluation cannot be added, narrow the title and claims to semantic similarity and repetition control rather than hallucination mitigation.

\subsection*{3. Hyperparameter selection appears to use the test set}

\classification{Likely test-set selection bias}

\findinglabel{Location.} Methodology; Experimental Setup; both result tables, especially ``Best Early-Stop'' and ``Runner ES.''

\findinglabel{Problem.} The method states $\alpha_0=18.0$ and $K=16$, but the sampling table uses $\alpha=15$ at 80 tokens and $\alpha=18$ at 200 tokens, while the greedy table reports both as ``Best'' and ``Runner'' configurations.  No validation set, search space, selection metric, or prespecified rule is described.  The labels therefore suggest that the best test-set result may have been selected after observing the same 500 test questions used for final reporting.

\findinglabel{Why it matters.} Selecting $K$, $\alpha_0$, schedule, decoding mode, or output budget on the test results biases effect estimates and invalidates nominal confidence intervals and $p$-values as confirmatory evidence.

\findinglabel{Suggested fix.} Provide disjoint train/vector-estimation, probe, validation, and test partitions.  State all partition sizes and group splits by source monograph and question before creating contrastive variants.  Select one configuration per decoding regime using validation data only, freeze it, and evaluate it once on the test set.  Move exploratory alternatives to an ablation table and remove ``Best'' terminology from test results.

\subsection*{4. Statistical comparisons and uncertainty reporting are not reproducible}

\classification{Definite reporting error and implementation ambiguity}

\findinglabel{Location.} Abstract; contribution~4; sampling-table caption and footnote; greedy table; Key Findings~1, 3, and 4.

\findinglabel{Problem.} The stars in the sampling table do not identify the reference condition for each test.  The caption promises 95\% confidence intervals, but the repetition column uses an undefined ``$\pm$'' quantity and latency has no interval.  The manuscript does not specify whether bootstrap tests are one- or two-sided, the bootstrap CI method, the paired statistic, or any multiple-comparison correction across methods, metrics, budgets, and decoding modes.  The greedy table contains no confidence intervals or tests even though contribution~4 says rigorous validation was conducted under both sampling and deterministic decoding.  The report $p=0.0000$ is invalid; with 10,000 resamples it should be reported as a bound such as $p<0.0001$ under an appropriate finite-resampling convention.

\findinglabel{Why it matters.} Readers cannot determine which differences are being tested, reproduce the analysis, or know whether nominal significance survives the large number of comparisons.  Ambiguous stars can attach a valid computation to an unsupported interpretation.

\findinglabel{Suggested fix.} For every comparison, state the reference, paired unit, statistic, alternative hypothesis, resampling method, number of resamples, CI construction, and multiplicity correction.  Define whether ``$\pm$'' denotes standard deviation, standard error, or a confidence interval.  Give paired effect estimates and 95\% CIs, not only $p$-values.  Apply the same reporting standard to greedy results or explicitly call them descriptive.  Replace $p=0.0000$ with a valid inequality.

\subsection*{5. Several headline numerical claims contradict the tables}

\classification{Definite quantitative errors}

\findinglabel{Location.} Main contribution~3; Key Findings~2 and 4; Conclusion.

\findinglabel{Problem.}
\begin{enumerate}
    \item Key Finding~2 calls BERTScore F1 $0.7189$ the highest value ``across the entire study,'' but the 80-token sampling table reports $0.7478$ and the 80-token greedy table reports $0.7477$.  At most, $0.7189$ is the highest value within the 200-token greedy condition.
    \item The claimed 99.97\% preservation is $34.32/34.33\approx99.97\%$, the ratio of two raw ROUGE-L scores.  It is not preservation of the steering \emph{gain}.  Relative to the greedy 80-token baseline, full steering gains $34.33-34.02=0.31$ points and early stopping gains $34.32-34.02=0.30$ points, so the retained ROUGE-L gain is about $96.8\%$.  For BERTScore in the same table, early stopping equals the baseline at $0.7467$, so it does not preserve the full-steering BERTScore gain.
    \item The Conclusion says repetition artifacts are ``eliminated,'' yet early stopping has higher 200-token repetition than vanilla under both sampling ($0.0372$ versus $0.0356$) and greedy decoding ($0.0385$ versus $0.0359$).  It is lower than full steering, not eliminated.
    \item Key Finding~4 writes ``$0.0572 - 0.0586$,'' which is typeset as subtraction rather than an interval and reverses the intended reading.
\end{enumerate}

\findinglabel{Why it matters.} These statements are central summary claims and will mislead readers even if the underlying table entries are correct.

\findinglabel{Suggested fix.} Restrict every superlative to its exact condition, calculate retained improvement relative to the baseline rather than dividing raw scores, replace ``eliminated'' with ``reduced relative to full steering,'' and write ``0.0572 and 0.0586'' or ``0.0572--0.0586.''  Recheck all abstract, contribution, finding, and conclusion numbers against a single analysis output.

\subsection*{6. Layer 8 and the truthfulness-vector interpretation are unsupported}

\classification{Unsupported claim and likely representation confound}

\findinglabel{Location.} Abstract; contributions; ``Contrastive Vector Estimation.''

\findinglabel{Problem.} The manuscript calls Layer~8 the ``optimal linear subspace,'' but reports no layer-wise probe design, metrics, held-out split, or comparison with other layers.  The notation says $h_l^+$ and $h_l^-$ are representations for responses $y_i^+$ and $y_i^-$, but Equation~(1) writes only $h_l^+(x_i)$ and $h_l^-(x_i)$; the response, token position, pooling rule, and chat formatting disappear from the mathematical definition.  It is therefore unclear whether the vector encodes truthfulness, synthetic perturbation templates, answer style, or final-token lexical artifacts.

\findinglabel{Why it matters.} Without held-out and source-disjoint probing, perfect or high separability can arise from dataset-construction artifacts.  The intervention may still affect outputs, but the paper cannot interpret the direction specifically as a transferable clinical-truthfulness representation.

\findinglabel{Suggested fix.} Add the complete layer-probing experiment, including every evaluated layer, question/source-disjoint splits, classifier and regularization, token position, class balance, and uncertainty.  Keep all positive/negative versions of one question in the same partition.  Write the representation explicitly as a function of both prompt and response, for example $h_l(x_i,y_i^{\pm})$, with a defined token index or pooling operator.  Test transfer to independently authored errors and prompt-only or pre-answer states.

\subsection*{7. The intervention is not specified sufficiently for reproduction}

\classification{Implementation ambiguity}

\findinglabel{Location.} Equations~(1)--(3); Experimental Setup.

\findinglabel{Problem.} The paper does not define whether $t=1$ is the prompt forward pass or the first generated token, whether the vector is added to all sequence positions or only the newest token, where Layer~8 is indexed and hooked, or how tuple outputs, KV caching, dtype, device placement, and hook removal are handled.  Equation~(3) gives $\alpha(K)=\alpha_0/K$, followed by an abrupt transition to zero at $K+1$; this is valid if intentional, but it does not decay to zero \emph{at} token $K$.  Full steering and random-direction construction are not mathematically defined.

\findinglabel{Why it matters.} Reasonable implementations can apply different numbers of interventions and produce different results, especially on the initial prompt pass and multi-GPU model shards.

\findinglabel{Suggested fix.} Provide pseudocode identifying the exact Hugging Face module, zero- or one-based layer convention, tensor slice (for example $H[:, -1, :]$), prompt-pass rule, counter update, cache behavior, dtype/device conversion, and random-vector normalization.  Clarify the endpoint of the linear schedule; if zero at $t=K$ is intended, use a denominator of $K-1$ for $K>1$.  Report software versions, batch size, prompt-length distribution, seeds, and decoding settings completely.

\subsection*{8. Dataset and annotation claims lack essential provenance}

\classification{Unsupported dataset-quality claim}

\findinglabel{Location.} Abstract; contribution~1; ``ViHaluEval-Medical Benchmark''; Experimental Setup.

\findinglabel{Problem.} The dataset is described only by size and three rounded category percentages.  The manuscript does not give the pharmacopoeia edition and citation, source-selection rules, question/answer generation procedure, perturbation rules, deduplication, split counts, licensing, or release status.  ``Expert-validated'' in the Abstract drifts to ``expert-annotated'' in the Introduction.  Cohen's $\kappa=0.91$ is uninterpretable without the number of raters, annotation unit, label set, double-annotation rate, prevalence, adjudication, and confidence interval.  If a panel has more than two raters, the use or aggregation of Cohen's two-rater statistic must be explained.

\findinglabel{Why it matters.} The claimed benchmark, layer selection, steering-vector estimation, and test validity all depend on leakage-free provenance and a clearly defined expert judgment process.

\findinglabel{Suggested fix.} Add a dataset card and complete split table.  Split by source monograph and question before generating contrastive variants, quantify residual similarity between partitions, and state explicitly that $\mathcal{D}_{steer}$ does not overlap validation or test sources.  Document expert qualifications and agreement computation, use one descriptor for the annotation process, cite the official source, and account for the percentages' rounded residual.

\subsection*{9. ``Zero latency overhead'' is contradicted by the reported measurements}

\classification{Definite wording error}

\findinglabel{Location.} Abstract; sampling and greedy tables; Conclusion.

\findinglabel{Problem.} At 80-token sampling, early stopping is reported at 10,746~ms versus 10,431~ms for vanilla, an increase of about 3.0\%, and it is 211~ms slower than full steering.  At 200-token greedy decoding, the best early-stop row is 23,020~ms versus 22,855~ms for vanilla, an increase of about 0.7\%.  Other rows fluctuate in the opposite direction.  No timing uncertainty, warm-up procedure, synchronization, batch size, or repetition count is given.

\findinglabel{Why it matters.} These measurements may indicate noise rather than a meaningful overhead, but they do not establish literal zero overhead.  Timing claims require a dedicated repeated-measures protocol.

\findinglabel{Suggested fix.} Report synchronized repeated timings after warm-up, including mean, dispersion, prompt and output lengths, batch size, hardware utilization, and the number of runs.  Use ``no measurable overhead under this protocol'' only if the confidence interval supports a practically negligible difference; otherwise report the observed overhead directly.

\section*{Equation, notation, sign, branch, and dimensional audit}

The direction $h_l^+-h_l^-$ followed by addition of $+\alpha(t)\hat v_{steer}^{(l)}$ is sign-consistent with the manuscript's intended truthful direction, and unit normalization makes the vector dimensionally compatible with the hidden state when $\alpha(t)$ is interpreted as an activation-scale coefficient.  No inverse-trigonometric functions, coordinate transformations, or branch/quadrant choices occur in the manuscript.  The concrete notation problems are the missing dependence on $y_i^{\pm}$ in Equation~(1), the undefined representation location, the schedule endpoint at $t=K$, the absence of mathematical definitions for full and random steering, and the drift between the fixed $\alpha_0=18$ in Methodology and $\alpha=15/18$ in the tables.  The term ``representation saturation'' is also used as a mechanism without any activation-norm, direction, or saturation measurement.

\section*{Meaningful editorial and reference findings}

\subsection*{The manuscript source does not compile}

\findinglabel{Location.} Related Work subsection heading ``Representation Steering \& ITI.''

\findinglabel{Problem.} A clean \texttt{pdflatex} run stops with \texttt{Misplaced alignment tab character \&} because the ampersand is unescaped in \verb|\subsection{Representation Steering & ITI}|.

\findinglabel{Why it matters.} The submitted source is not buildable as written, so later layout and reference warnings cannot be validated from a clean build.

\findinglabel{Suggested fix.} Change the heading to \verb|\subsection{Representation Steering \& ITI}|, rebuild at least twice, and inspect the final log for unresolved citations, overfull boxes, and table-footnote layout.

\subsection*{Citation coverage and relevance are insufficient}

\findinglabel{Location.} Introduction; Related Work; benchmark description; metrics and statistical analysis; bibliography.

\findinglabel{Problem.} The official pharmacopoeia, BERTScore, ROUGE-L implementation, repetition metric, and statistical procedures are uncited.  Citation key \texttt{ref1} is used to support SFT, DPO, and catastrophic forgetting even though the listed item is a representation-engineering paper, so the citation does not directly support the sentence.  The assertion that standard ITI causes long-output fluency degradation also needs direct evidence or a citation.

\findinglabel{Why it matters.} Citation mismatch obscures which claims are established by prior work and weakens the novelty and factual basis of the paper.

\findinglabel{Suggested fix.} Retain \texttt{ref1} and \texttt{ref2} for their actual representation-intervention contributions, add appropriate sources for SFT/DPO and forgetting, cite the pharmacopoeia directly, and add the original metric references and relevant multilingual/medical validation literature.  Cite or experimentally substantiate the continuous-steering degradation claim.

\section*{Proofing sweep}

The bundled mechanical scan was run on both the current source and the available PDF.  The source produced no high-confidence pattern hits.  The PDF produced only a duplicated-period hit from the mathematical ellipsis in $t\in\{1,2,\ldots,N\}$; this is intentional and is not a punctuation error.  The bibliography was checked for duplicated commas, malformed quotation punctuation, and obvious title-format defects; no such defect was found.

The following bounded manual corrections are high-confidence:
\begin{itemize}
    \item \textbf{Key Finding~4:} replace the mathematical subtraction ``$0.0572 - 0.0586$'' with ``0.0572 and 0.0586'' or the typographic range ``0.0572--0.0586.''
    \item \textbf{Key Finding~4:} replace $p=0.0000$ with a valid lower bound determined by the resampling procedure, such as $p<0.0001$ if justified.
    \item \textbf{Author block:} ``Kien'' and ``Corresponding Authors: Phan Do Thanh Tuan, Kien.'' appear incomplete; provide each author's full publication name and a consistent corresponding-author designation.
    \item \textbf{Terminology:} use ``early stopping'' as a noun and ``early-stopping'' as an attributive adjective; standardize capitalization of ``sign-flipped'' and ``full steering'' in running prose.
    \item \textbf{Tables:} define what boldface means.  Some bold entries mark the desired minimum, while the sign-flipped repetition maximum is also bold, making the visual convention ambiguous.
\end{itemize}

\section*{Items requiring external verification}

\begin{itemize}
    \item Verify and cite the exact official title, edition, issuing authority, date, and reuse permissions of the Vietnamese pharmacopoeia used to build the dataset.
    \item Verify the clinical correctness of all dosage, pregnancy/lactation, and interaction reference answers against current authoritative sources; disclose whether experts reviewed every record or a sampled subset.
    \item Validate the chosen BERTScore model and configuration for Vietnamese clinical semantics, including model identifier, tokenizer, IDF use, baseline rescaling, and evidence that score differences correspond to clinically meaningful changes.
    \item Verify the characterization of Qwen2.5-7B-Instruct as an SLM and the untested claims of privacy preservation and local-hospital deployability against an explicit operational definition and deployment measurements.
    \item Verify novelty against recent work on activation steering schedules, multilingual medical hallucination evaluation, and inference-time intervention; the current five references are insufficient for a novelty claim.
\end{itemize}

\section*{Highest-priority revision sequence}

\begin{enumerate}
    \item Separate schedule effects from steering-strength effects using equal-$\alpha_0$ and cumulative-dose-matched ablations.
    \item Add a clinically adjudicated hallucination/safety endpoint; do not infer medical correctness from BERTScore or repetition alone.
    \item Establish disjoint validation and test protocols, freeze hyperparameters before test evaluation, and rerun the analysis.
    \item Correct the statistical specification, comparison targets, confidence intervals, multiplicity handling, and invalid $p$-value reporting.
    \item Correct the contradictory numerical claims, especially ``highest across the entire study,'' 99.97\% gain retention, ``eliminated,'' and ``zero latency overhead.''
    \item Add complete layer-probing, vector-construction, hook, data, annotation, and environment details for reproduction.
    \item Repair the LaTeX compilation error, complete the author metadata, and expand and correct the citations.
\end{enumerate}

\end{document}